\scanpage{46}
\uline{Ens de card 2} : torseur sur \(\mathbb F_2\) i.e. \deleted{\unsure{espace}} \inserted{droite} affine \deleted{\unsure{sur une}} \(\mathbb F_2\)

\uline{Ens de card 3} : droite affine sur \(\mathbb F_3\) / éléments non nuls d'un vectoriel de dim \(2\) sur \(\mathbb F_2\) / droite projective sur \(\mathbb F_2\) / \(+\) orient\textsuperscript{n} : \deleted{\unsure{espace affine} \(\mathbb F_3\)}

\uline{Ens de card 4} : \deleted{\unsure{droite}} plan affine sur \(\mathbb F_2\) / droite projective sur \(\mathbb F_3\) / \deleted{\unsure{sens}} \uline{orient\textsuperscript{n}} : droite affine sur \(\mathbb F_4\) ?

\uline{Ens de card 5} \((\Leftrightarrow\) droite projective sur \(\mathbb F_5\), les \(6\) pts sont les sommets du \unsure{biicosaèdre|biicosoèdre} associé à \(E)\)

\uline{Ens orienté de card 5} droite projective sur \(\mathbb F_4\) -- \unsure{l'énon.|l'énum.|l'ens.} :

Cas \(G=\mathrm{GP}(2,\mathbb F_q)\) : \inserted{\(r=2\)} \(N=3,\) \deleted{\unsure{\(r\)}} \(d_\lambda=2,3\)

\[\boxed{\operatorname{card}G(\mathbb F_q):q^N(q-1)^r\prod_{1\leq\lambda\leq r}\frac{\overbrace{q^{d_\lambda}-1}^{\displaystyle1+q+\cdots+q^{d_\lambda-1}}}{q-1}}\]
\[P_{G/B}=\frac{(1-T^2)^r}{\prod_\lambda(1-T^{2d_\lambda})}\]

\[\mathrm{SL}(3,\mathbb F_q)\quad\operatorname{card}\ \sum_i\underline b\,q^{d(i)}\qquad b=q^3(q-1)^2\sum_{0\leq i\leq3}b_{2i}q^i\]
\[(1:P_{G/B}(\sqrt q))=\frac{(1-q)^r}{\prod_\lambda(1-q^{d_\lambda})}=\frac{(q-1)^r}{\prod_\lambda(q^{d_\lambda}-1)}\]

\[\operatorname{card}G(\mathbb F_q)=\deleted{\unsure{\lvert G\rvert}}\ q^N(q-1)^r\underbrace{\operatorname{card}G/B(\mathbb F_q)}_{\displaystyle\sum_{\alpha\text{ cellule}}q^{d(\alpha)}}\]
\[\begin{aligned}{}&=\sum_{0\leq i\leq N}b_{2i}q^i=P_{G/B}(\sqrt q)\ \Bigg\{\\[-2pt]&(d(\alpha)\ \text{sa dim})\end{aligned}\]
\margin{gr. de Chevalley \quad \deleted{\unsure{cellules}}}
\[\begin{aligned}N&=\text{nb de racines pos.}\\&=\dim G/B=\dim B\\r&=\operatorname{rg}G\end{aligned}\]
\diagram{\(\begin{array}{c|r|r|r|r|l}q&2&3&4&5&\\ \hline &1&2&3&4&\mathrm{GL}(\unsure{1|2},\mathbb F_q)\deleted{\unsure{\,2\,4}}\\ \cline{2-5}&2&6&12&20&\mathrm{Aff}(1,\mathbb F_q)\\ \cline{2-6}&6&24&60&120&\mathrm{SL}(2,\mathbb F_q)\\ \cline{2-6}&6&24&60&120&\mathrm{GP}(1,\mathbb F_q)\\ \cline{2-6}&6&48&180&480&\mathrm{GL}(2,\mathbb F_q)\end{array}\)}
\[\unsure{1,2}\qquad 24\mid9.48\mid\underline{16.180}\mid25.480\mid\quad\mathrm{Aff}(2,\mathbb F_q)\]
\diagram{\begin{tikzcd}[ampersand replacement=\&,column sep=small,row sep=small] \unsure{B_T|B\,T}\arrow[d,dash] \&[-4pt] \supset G/T\simeq G/B \\ B_G \& P_{B_G}\cdot P_{G/B}=P_{B_T}\\ \& \displaystyle\left(\frac1{1-T^2}\right)^r\ \overset{\text{\unsure{d'où}}}{\unsure{\longleftarrow|\longleftrightarrow}}\ \prod_\lambda\frac1{1-T^{2d_\lambda}}\end{tikzcd}}
\[\begin{aligned}\mathrm{GL}(1,\mathbb F_q)&\simeq\mathbb F_q^*\simeq\mathbb Z/(q-1)&\operatorname{card}\ q-1\\ \mathrm{Aff}(1,\mathbb F_q)&\simeq\mathbb F_q^*\cdot\mathbb F_q&\operatorname{card}\ q(q-1)\\ \mathrm{GP}(1,\mathbb F_q)&\simeq\mathrm{GL}(2,\mathbb F_q)/\mathbb F_q^*\end{aligned}\]
\[\begin{aligned}\inserted{\text{card}}\ \mathrm{SL}(2,\mathbb F_q)&=\sum_{\text{cell. }i}b\cdot q^{d(i)}\qquad b=q(q-1)\\ \operatorname{card}\mathrm{GP}(1,\mathbb F_q)&=(\unsure{q+1})q(q-1)\\ \operatorname{card}\mathrm{GL}(2,\mathbb F_q)&=q(q-1)^2(q+1)\end{aligned}\]
\[\begin{aligned}\operatorname{card}\mathrm{SL}(1,\mathbb F_q)&=q(q-1)(q+1)\\ \operatorname{card}\mathrm{Aff}(2,\mathbb F_q)&=q^3(q-1)^2(q+1)\end{aligned}\]
% archive sheet number: 45
